\chapter{Autenticazione su Internet}

\section{Kerberos}

Kerberos è uno standard Internet e lo standard \textit{de facto} per l'autenticazione remota in reti locali aziendali, sviluppato originariamente al Massachusetts Institute of Technology (MIT). 

\begin{notebox}[Il Modello Kerberos]
È un servizio di autenticazione basato su una \textbf{terza parte fidata} (\textit{trusted third party}). Client e server non si autenticano scambiandosi direttamente password, ma si affidano a Kerberos per mediare la loro autenticazione reciproca, richiedendo all'utente di provare la propria identità per ogni servizio invocato in modo sicuro.
\end{notebox}

\subsection{Il Protocollo Kerberos}

Il protocollo Kerberos coinvolge tre entità principali: il client, i server applicativi e il server centrale Kerberos (composto a sua volta da due servizi). È progettato per contrastare varie minacce di rete (in particolare gli attacchi di impersonificazione e \textit{replay}) centralizzando l'autenticazione in modo sicuro tramite crittografia simmetrica (originariamente DES, ora tipicamente AES).

L'\textbf{Authentication Server (AS)} conosce tutte le password (hash) di tutti i client registrati. Non si limita a validare l'identità, ma fornisce al client un \textit{ticket} crittografico per richiedere servizi in modo sicuro, senza mai far viaggiare la password in chiaro.

\subsubsection{Funzionamento del Protocollo}
Il flusso di autenticazione si articola in tre fasi e sei passaggi ben definiti:

\begin{enumerate}
    \item \textbf{Richiesta di un Ticket-Granting Ticket (TGT):}
    L'utente si collega a una workstation (client) e invia una richiesta iniziale all'AS contenente semplicemente il proprio ID utente in chiaro. 
    L'AS verifica l'ID, crea un Ticket-Granting Ticket (TGT) e una chiave di sessione temporanea. Entrambi vengono \textbf{cifrati con una chiave derivata dalla password dell'utente} e inviati indietro al client.
    
    \item \textbf{Ottenimento del TGT da parte del Client:}
    Il client richiede la password all'utente, calcola la chiave derivata e decifra il messaggio dell'AS, recuperando il TGT e la chiave di sessione. Il TGT è una credenziale riutilizzabile (cifrata con una chiave segreta condivisa esclusivamente tra AS e il Ticket-Granting Server - TGS) contenente l'ID utente, un timestamp, una scadenza temporale e una copia della chiave di sessione.
    
    \item \textbf{Richiesta di un Service-Granting Ticket (SGT):}
    Per accedere a uno specifico server applicativo (es. Server V), il client invia al TGS una richiesta contenente il TGT ottenuto precedentemente e un \textit{Authenticator}.
    
    \begin{defbox}[L'Authenticator di Kerberos]
    Un \textbf{Authenticator} è una struttura dati crittografica \textit{monouso} (one-time) e a vita molto breve (pochi minuti), contenente l'ID del client e il timestamp attuale, cifrata con la chiave di sessione. Poiché scade velocemente, previene gli attacchi di \textit{replay} in cui un attaccante tenta di riutilizzare un ticket intercettato.
    \end{defbox}
    
    \item \textbf{Concessione del Service-Granting Ticket:}
    Il TGS, essendo la stessa entità di Kerberos, decifra il TGT (in cui trova la chiave di sessione) e usa quest'ultima per verificare l'Authenticator. Se valido, crea un Service-Granting Ticket (SGT) specifico per il Server V e lo invia al client insieme a una \textit{nuova chiave di sessione specifica per il servizio}, il tutto cifrato con la chiave di sessione originale. Il SGT è cifrato con una chiave segreta condivisa solo tra TGS e il Server V.
    
    \item \textbf{Richiesta del Servizio al Server:}
    Il client inoltra al Server V il SGT appena ricevuto e un nuovo Authenticator. Il Server V decifra il ticket con la propria chiave privata, verifica l'Authenticator e, se le credenziali corrispondono, concede l'accesso al servizio. 
    
    \item \textbf{Autenticazione Mutua (Opzionale ma consigliata):}
    Il server risponde con un messaggio cifrato (es. il timestamp dell'Authenticator incrementato di 1) per dimostrare al client la propria autenticità, impedendo il \textit{server spoofing}.
\end{enumerate}

\begin{center}
\begin{tikzpicture}[
    node distance=2.5cm and 4.5cm,
    box/.style={rectangle, draw=blue!80!black, thick, minimum width=2.8cm, minimum height=1.2cm, align=center, fill=blue!5, rounded corners},
    arrow/.style={-{Stealth}, thick, font=\scriptsize, blue!80!black}
]

\node[box] (client) {\textbf{Client}\\(Utente)};
\node[box, fill=purple!10, above right=of client] (as) {\textbf{AS}\\(Auth Server)};
\node[box, fill=orange!10, right=of client] (tgs) {\textbf{TGS}\\(Ticket-Granting Server)};
\node[box, fill=green!10, below right=of client] (server) {\textbf{Server V}\\(Applicativo)};

% Linee verso AS
\draw[arrow] (client) -- node[above left, align=right] {1. Richiesta TGT\\(ID Utente)} (as);
\draw[arrow, red] ([yshift=-2mm]as.south west) -- node[below right, align=left] {2. TGT + Session Key\\(Cifrati con Password)} ([xshift=2mm]client.north east);

% Linee verso TGS
\draw[arrow] ([yshift=2mm]client.east) -- node[above] {3. TGT + ID Servizio + Authenticator} ([yshift=2mm]tgs.west);
\draw[arrow, red] ([yshift=-2mm]tgs.west) -- node[below] {4. SGT + Nuova Chiave Servizio} ([yshift=-2mm]client.east);

% Linee verso Server
\draw[arrow] (client) -- node[below left, align=right] {5. SGT +\\Nuovo Authenticator} (server);
\draw[arrow, red] ([yshift=2mm]server.north west) -- node[above right, align=left] {6. Mutua Autenticazione\\(Opzionale)} ([xshift=2mm]client.south east);

\end{tikzpicture}
\captionof*{figure}{Flusso di autenticazione del Protocollo Kerberos}
\end{center}

\subsection{Realm Kerberos e Kerberi Multipli}
Un ambiente Kerberos completo e isolato (composto da un server Kerberos principale, vari client e server applicativi) è definito \textbf{realm}, e coincide generalmente con un dominio amministrativo di rete. 

Per consentire l'autenticazione \textit{inter-realm} (cioè un utente di un realm che accede a risorse di un altro), i server Kerberos di due realm diversi devono condividere una chiave segreta stabilita amministrativamente e fidarsi reciprocamente, permettendo così al client di ottenere un TGT locale per richiedere un ticket a un TGS remoto.

\subsection{Versioni di Kerberos}
\begin{itemize}
    \item \textbf{Versione 4:} La prima ampiamente utilizzata, ma presentava vincoli di sicurezza (basata esclusivamente sul cifrario DES).
    \item \textbf{Versione 5:} Introdotta nel 1993 (e integrata nativamente in \textit{Microsoft Active Directory}). Introduce migliorie sostanziali, tra cui il supporto per algoritmi crittografici multipli moderni (come AES), l'\textit{Authentication Forwarding} (delega delle credenziali a un server intermedio per agire per conto del client) e una maggiore scalabilità per l'autenticazione inter-realm (cross-realm trust).
\end{itemize}

\section{X.509}

Lo standard X.509 (definito dall'ITU-T) è il formato universalmente accettato per i certificati a chiave pubblica, utilizzati trasversalmente in tecnologie come IPsec, TLS/HTTPS, S/MIME, VPN e firme elettroniche. 

\begin{defbox}[Cos'è un Certificato X.509?]
Un certificato lega crittograficamente una \textbf{chiave pubblica} all'\textbf{identità} del suo proprietario (soggetto). Per garantirne la validità, il certificato è firmato digitalmente da una terza parte universalmente fidata, chiamata \textbf{Certificate Authority (CA)}.
\end{defbox}

\subsection{Formato del Certificato}
Un certificato X.509 (versione 3) include elementi fondamentali strutturati in modo rigoroso:
\begin{itemize}
    \item \textbf{Nome del Soggetto:} L'entità (persona, server web o organizzazione) proprietaria della chiave.
    \item \textbf{Chiave Pubblica del Soggetto:} Il valore effettivo della chiave e l'algoritmo (es. RSA 2048-bit).
    \item \textbf{Periodo di validità:} Date di "Not Before" e "Not After".
    \item \textbf{Nome dell'Emittente:} L'identità della CA che ha rilasciato il certificato.
    \item \textbf{Firma Digitale della CA:} L'hash crittografico di tutto il certificato, firmato con la chiave privata dell'emittente, che garantisce l'inviolabilità dei dati.
\end{itemize}

La \textbf{versione 3} dello standard ha introdotto estensioni flessibili e critiche, come i \textit{Basic Constraints}, i quali specificano se un certificato appartiene a una CA (ed è quindi abilitato a firmare a sua volta altri certificati) oppure appartiene a un utente finale (end-entity).

\subsection{Revoca dei Certificati}
Un certificato può essere revocato prima della sua naturale scadenza temporale (ad esempio in caso di furto o compromissione della chiave privata del soggetto). I meccanismi principali sono due:
\begin{description}
    \item[Certificate Revocation List (CRL):] Liste pubblicate periodicamente e firmate dalla CA, contenenti i numeri di serie dei certificati non più validi. La loro gestione su larga scala è onerosa e poco reattiva.
    \item[Online Certificate Status Protocol (OCSP):] Un'alternativa moderna e più leggera che permette a un'applicazione (es. un browser web) di interrogare direttamente e dinamicamente i server della CA sullo stato in tempo reale di un singolo certificato specifico.
\end{description}

\section{Public-Key Infrastructure (PKI)}

Una PKI (Public-Key Infrastructure) è l'insieme complesso di hardware, software, persone, policy aziendali e procedure amministrative necessarie per creare, gestire, distribuire, utilizzare, conservare e revocare in sicurezza i certificati digitali su vasta scala.

\subsection{Modello di Trust e Problematiche Attuali}
La verifica autonoma di un certificato richiede di fidarsi incondizionatamente della CA emittente, creando una \textbf{"catena di fiducia" (Trust Chain)} che parte dal certificato del server, passa per CA intermedie, fino a risalire a una \textit{Root CA} (CA Radice) fondamentale. 

I sistemi operativi e i browser web attuali includono un \textit{trust store} (archivio di certificati fidati) contenente decine o centinaia di Root CA preinstallate. Questo modello ha dimostrato diverse criticità pratiche:
\begin{warnbox}[Debolezze della PKI Pubblica]
\begin{itemize}
    \item \textbf{Fiducia orizzontale incondizionata:} Il sistema presume che \textit{tutte} le CA nel trust store siano ugualmente affidabili e competenti. La compromissione di una singola CA debole (come nel disastroso incidente di DigiNotar del 2011) permette di generare certificati falsi ma tecnicamente validi per qualsiasi dominio del mondo.
    \item \textbf{Decisioni dell'utente:} Gli utenti finali spesso non hanno le competenze tecniche per valutare adeguatamente gli avvisi di sicurezza generati dai browser in caso di certificati non validi, portando a bypassarli.
    \item \textbf{Trust store frammentati:} Browser diversi (es. Firefox vs Chrome) e Sistemi Operativi (es. Windows vs macOS) gestiscono e aggiornano trust store del tutto differenti, portando a visioni di sicurezza non uniformi.
\end{itemize}
\end{warnbox}

\subsection{Public Key Infrastructure X.509 (PKIX)}
Per cercare di standardizzare queste dinamiche, il gruppo di lavoro \textbf{PKIX} della IETF ha definito un modello architetturale formale per l'implementazione di una PKI su Internet. Il modello identifica rigorosamente gli elementi chiave del sistema (\textit{End entity}, \textit{CA}, \textit{Registration Authority} o RA, e \textit{Repository} CRL) e le relative funzioni di gestione (registrazione degli utenti, certificazione, recupero chiavi e revoca).
